\section{Why one acceptable task per person is not enough}
Two students, Ada and Ben, each accept task $1$. Neither accepts any other task. Individually both have an acceptable choice. Together they cannot receive distinct tasks. The problem is collective competition for the same resource.

Now suppose Ada accepts tasks $1,2$, while Ben accepts only $1$. A solution exists: Ada receives $2$ and Ben receives $1$. But a greedy procedure that gives Ada task $1$ immediately gets stuck. The failure of that procedure is not a proof of impossibility. It only says the current assignment cannot be extended without revising an earlier choice.

A bipartite graph models the problem. Its vertices are split into disjoint sets $L$ and $R$, and every edge joins one vertex in $L$ to one in $R$. Let $L$ be students, $R$ tasks, and an edge indicate acceptability. The word bipartite means that edges never connect two students or two tasks in this model.

A matching is a collection of edges with no shared endpoints. It assigns distinct tasks to the students that it covers. A matching \emph{saturates} $L$ when every student is covered. It need not use every task. This distinction matters when there are more tasks than students.
